\documentclass[00-main.tex]{subfiles}

\usepackage{parskip} % no indents

\newcommand{\reviewercomment}[1]{\vspace{0.3cm}\noindent\textbf{#1}\vspace{0.15cm}}
\newcommand{\manuscriptquote}[1]{\begin{quote}\textit{#1}\end{quote}}

\begin{document}

Dear \emph{Neural Computation} Editor and Reviewer,

\vspace{0.5cm}

We sincerely thank the reviewer for their excellent and thoughtful feedback on our manuscript, ``Quantifying uncertainty in drift diffusion models of decision making
under temporal dependence and parameter variability,'' and appreciate the time and care invested in raising .... We hope our ... is reflected in the revisions and additions we have undertaken. We have revised the text to address the reviewer's comments. As an overview, we have:

\begin{enumerate}
    \item Summarize changes
    \item ...
\end{enumerate}

\vspace{0.5cm}

Below, we provide a point-by-point response to the reviewer's comments. The reviews are in \textbf{bold}, our responses in plain text, and changes to the manuscript in indented \textit{italics}. Please also see the revised manuscript (clean and tracked versions).

\subsubsection*{Reviewers' comments}

\reviewercomment{1. The authors identify four limitations of current inference methods and present pseudo-MLE as the remedy. However, Bayesian MCMC sampling already addresses all four of these issues, at least in principle:}

\reviewercomment{- Uncertainty quantification: MCMC provides full posterior distributions;}

Response ...

\manuscriptquote{Manuscript quote ...}

\reviewercomment{- Asymptotic approximations: This seems to be an issue of model misspecification (i.e., the lack of temporal dependencies in the model), which applies to any estimation method as far as I know;}

Response ...

\manuscriptquote{Manuscript quote ...}

\reviewercomment{- Computational cost: The authors' own Figure 7A shows that the worst-case MCMC runtime in their comparison was approximately 15 minutes. This doesn't seem like a meaningful bottleneck. The computational advantage of pseudo-MLE would be more compelling in settings with very large datasets;}

Response ...

\manuscriptquote{Manuscript quote ...}


\reviewercomment{- Robustness to misspecification and temporal correlattions: This is the most substantive claim. The authors show in Figure 4 that sandwich estimators outperform non-robust estimators under misspecification. However, given that the data-generating process includes trial-by-trial variability in t0, v, and z, it would make sense to compare their pseudo-ML method to fitting the full 7-parameter DDM that includes these trial-by-trial variabilities. If the seven-parameter DDM recovers parameters and uncertainties accurately in this setting, the case for the sandwich correction is considerably weakened. Furthermore, the comparison in Figure 7 is conducted under correct specification and iid conditions. A more informative comparison would evaluate MCMC versus pseudo-MLE under the misspecified, temporally dependent conditions that motivate the paper. Without this, the robustness advantage of pseudo-MLE over MCMC remains undemonstrated.}

Response ...

\manuscriptquote{Manuscript quote ...}

\reviewercomment{2. The proposed method is restricted to the four-parameter DDM. This is a notable constraint. The four-parameter DDM is known to systematically misfit the relative speed of error responses observed in human reaction time data (e.g., Ratcliff \& McKoon, 2008; Ratcliff et al., 2016). This motivated the development of the 6-parameter (and later 7-parameter) DDM in the first place (Ratcliff, 1978). Adopting the proposed method therefore requires accepting a model that is known to be empirically inadequate for a large class of paradigms (at least those involving human participants).}

\reviewercomment{The authors acknowledge that MCMC-based approaches (e.g., HSSM) can accommodate approximate likelihoods and thus extend to the seven-parameter DDM, collapsing-bound models, and other variants. This flexibility is central to whether a method is practically useful across the range of experimental designs in the field. The authors should more explicitly discuss the conditions under which the four-parameter restriction is acceptable, and for which research questions the proposed method is and is not appropriate.}

Response ...

\manuscriptquote{Manuscript quote ...}

\reviewercomment{3. Empirical novelty. Many of the ideas put forward - modelling parameters as a function of covariates - have already been described and applied succesfully in the literature, including in combination with MCMC sampling to estimate uncertainty. E.g., the cited Cochrane et al. (2023) paper used time-on-task as covariates; the Gunawan et al. (2020) paper used both Hidden Markov models and time on task as covariates; the Kucharský et al. (2021) paper also used Hidden Markov models; Alister \& Evans (2025) use time on task; Miletić et al. (2025) do something similar, and also use delta rules borrowed from reinforcement learning. The latter is often done in the RL-DDM tradition (e.g., Pedersen et al., 2017; Fontanesi et al., 2019), which also uses covariates (in that case, reward) to define latent trajectories of trial-by-trial parameters.}

\reviewercomment{The authors cite some of these works but do not adequately differentiate their contribution from them. The paper would benefit from a clearer explanation of what the pseudo-MLE framework offers that these approaches do not.}

Response ...

\manuscriptquote{Manuscript quote ...}

\end{document}